\documentclass{article}
\usepackage{amsmath,amsthm}
\newcommand{\R}{\mathbb{R}}
\newcommand{\pair}[2][x]{#1 + #2}
\newtheorem{theorem}{Theorem}
\title{Native LaTeX reading}
\author{Example document}
\date{}

\begin{document}
\maketitle

\section{Reading and editing}
Choose Source in one pane and Rendered in another. Edit this paragraph:
the rendered pane updates before you save. Text can be \textbf{bold},
\emph{emphasized}, or \texttt{monospaced}. Inline math uses the document's
macros: $f : \R \to \R$ and $\pair{y}$.

\begin{quote}
Rendered mode reflows to the pane. PDF Preview remains available for exact
typesetting of saved files.
\end{quote}

\section{Equations}
\begin{align}
  (x+y)^2 &= x^2 + 2xy + y^2 \label{square} \\
  x^2-y^2 &= (x-y)(x+y) \label{difference}
\end{align}
Equation \eqref{square} and equation \eqref{difference} keep their numbers
when you scroll or change panes.

\begin{theorem}[An example]
For every $x \in \R$, $x^2 \geq 0$.
\end{theorem}
\begin{proof}
This text and its math are laid out natively.
\end{proof}

\section{Lists and tables}
\begin{enumerate}
  \item Source edits the shared document.
  \item Rendered follows unsaved changes.
  \item PDF Preview compiles saved files.
\end{enumerate}

\begin{tabular}{lr}
  Mode & Compiler required \\ \hline
  Source & No \\
  Rendered & No \\
  PDF Preview & Yes \\
\end{tabular}

\begin{verbatim}
Commands inside verbatim stay literal: \section{Example} and $x^2$.
\end{verbatim}
\end{document}
